% Generated draft: requires mathematical and source review.
\documentclass[11pt]{article}
\usepackage[margin=25mm]{geometry}
\usepackage{fontspec}
\setmainfont{Latin Modern Roman}
\usepackage{amsmath,amssymb,mathtools,unicode-math}
\setmathfont{Latin Modern Math}
\usepackage{enumitem,needspace,xurl,hyperref}
\setlength{\parindent}{0pt}
\setlength{\parskip}{5pt}
\setlength{\emergencystretch}{3em}
\newcommand{\N}{\mathbb{N}}
\newcommand{\Z}{\mathbb{Z}}
\newcommand{\Q}{\mathbb{Q}}
\newcommand{\R}{\mathbb{R}}
\newcommand{\C}{\mathbb{C}}
\newcommand{\F}{\mathbb{F}}
\newcommand{\A}{\mathbb{A}}
\newcommand{\PP}{\mathbb{P}}
\newcommand{\E}{\mathbb{E}}
\newcommand{\eps}{\varepsilon}
\newcommand{\dd}{\,\mathrm{d}}
\begin{document}
\section*{UnsolvedMath ID: 10400118}
\textbf{Source problem code:} catalogue entry
\subsection{Definitions and mathematical statement}
A rational homology \(3\)-sphere is a closed oriented \(3\)-manifold \(M\) with \(H_i(M;\Q)\cong H_i(S^3;\Q)\) for every \(i\). Let \(G\) range over all simply connected compact simple Lie groups, and for each \(G\) let \(r\) range over all levels allowed for the corresponding quantum-invariant construction. Denote the source's quantum \(G\) and quantum \(PG\) invariants by \(\tau_r^G(M)\) and \(\tau_r^{PG}(M)\), respectively. The definition of the \(PG\) variant, allowed levels, root choices and normalizations are not given in this extract.

S. K. Hansen and T. Takata's Problem 7.3 asks whether each of these families separates rational homology \(3\)-spheres. Under the usual oriented-homeomorphism meaning of separation, the first question is whether, for all rational homology \(3\)-spheres \(M,N\),
\[
 \left(\forall G\ \forall r\text{ allowed for }G:\quad
 \tau_r^G(M)=\tau_r^G(N)\right)
 \Longrightarrow M\cong_{\rm or}N.
\]
The analogous question replaces \(\tau_r^G\) by \(\tau_r^{PG}\), with the allowed levels for that construction. These are questions about each full family, not merely the union of both families or a fixed group or level.

The problem also asks how well these families separate arbitrary closed oriented \(3\)-manifolds. In particular it seeks an account of the indistinguishability relation given by equality of all invariants in a chosen family; the source supplies no narrower quantitative criterion. Its precise manifold equivalence convention and all construction parameters must be recovered before an unconditional separation statement can be accepted.
\subsection{Short English statement}
Do all allowed-level quantum invariants, ranging over every simply connected compact simple Lie group, determine a rational homology three-sphere? Ask this separately for the \(G\) and \(PG\) families, then investigate their distinguishing power for general closed oriented three-manifolds.
\subsection{Sources}
\begin{itemize}
\item Local UnsolvedMath registry, permanent ID \(10400118\), exact source code \texttt{catalogue entry}, Problem 7.3, attributed to S. K. Hansen and T. Takata. The supplied title is truncated after ``(S.K.''
\item Ohtsuki, \emph{Problems on invariants of knots and 3-manifolds} (2002), Problem 7.3, PDF page 101: \url{https://msp.org/gtm/2002/04/gtm-2002-04-024s.pdf}. Bold headings distinguish 231 numbered items from surrounding references.
\item Extraction pdf-font-structure, accessed 2026-07-29. Source evidence NEEDS\_REVIEW presents an open problem and requires a current primary-literature status audit. Supplied OPEN-TRIAGE summary says separation of rational homology spheres is not obviously resolved. No further primary reference or parameter definition is supplied.
\item Dataset attribution: Ulam AI / UnsolvedMath Contributors, CC BY 4.0. Source rights NEEDS\_REVIEW: publicly accessible author/publisher copy, redistribution terms requiring review.
\end{itemize}
\end{document}
